\section{Gaussian scalar extensions of the Weyl groups \(E_6,E_7\)}
\label{sec:e6-e7-gaussian}

\subsection*{D-E67-01. Exact local data from the published class representatives}

Let \(K=\mathbb Q(i)\), and let \(W_m=W(E_m)\), for \(m=6,7\),
act in its rational simple-root basis. Its invariant positive form is
the Cartan matrix \(H_m\): diagonal entries are two, and its edges
are
\[
 (1,3),(3,4),(2,4),(4,5),(5,6),
 \quad\text{and additionally }(6,7)\text{ for }m=7.
\]
Thus \(\det H_6=3\), \(\det H_7=2\), and the simple reflections
are explicitly
\[
                   s_j=I-e_j(H_m)_{j,*}.
\]
We consider \(J_m=\mu_4W_m\). All statements in the following finite
calculation use exact rational Gaussian arithmetic.

Carter's published classification gives 25 classes for \(W_6\) and
60 for \(W_7\), with their names and sizes in
\cite[Table 9, p.~50; Table 10, pp.~51--52]{WG:Carter}.
The author's CHEVIE library supplies words in the simple reflections for
these classes. The precise data are \texttt{ClassInfo.classtext} in
\texttt{weyle6.g}, and \texttt{WordsClassRepresentatives} in
\texttt{weyle7.g}; CHEVIE's manual, \S87.1, documents these as
complete conjugacy-class representatives. We use this published
classification and documented data for completeness. We do not claim
an independent enumeration of all 2,903,040 elements of \(W_7\).

The compact retained input files contain only the Cartan matrices, class
names, words, orders and sizes, together with the URL and SHA256 of each
original source. The checker reconstructs each matrix from the displayed
simple reflections, verifies its exact order and orthogonality, and
computes each of \(w,iw,-w,-iw\). For each such element \(g\), it
takes independent columns of \(1-g\) as a basis of \(U_g\), and
computes
\[
       d_g=\det(H_m|_{U_g}),\qquad
       b_g=\det(1-g|_{U_g}),\qquad k_g=\dim_K U_g.
\]
Finite Fourier inversion modulo the appropriate cyclotomic polynomial
gives exact eigenvalue multiplicities and ages.

In the rank-seven case, \(b_g\) may be a nonzero rational multiple
of \(1+i\) or \(1-i\), as well as a real or imaginary rational
number. Write
\[
              |b_g|=c_g\,2^{\epsilon_g/2},
           \qquad c_g\in\mathbb Q_{>0},\quad\epsilon_g\in\{0,1\},
           \qquad q_g=d_g/c_g.
\]
This gives a rational local norm class. In a product-one tuple of even
total moving rank, \(\prod b_g\) is rational by D-DET-01. Its phase
then forces \(\sum\epsilon_g\) to be even. The removed factor
\(2^{\sum\epsilon_g/2}\) is a Gaussian norm. For a balanced tuple
the sign of the product is positive, also by D-DET-01. Therefore the
signed determinant class of the parabolic space is precisely
\[
                              \left[\prod q_g\right].
\]
Changing moving-space bases multiplies each \(q_g\) by a Gaussian norm.

The complete calculations give the following particularly small table.
Here ``nonreal'' means \(g\in i\mathbb R^{m\times m}\).
\[
\begin{array}{c|c|c|c|c}
m&\min k_g\text{ for nonreal }g&
 \text{underlying Carter classes attaining it}&q_g&
 \text{age}(iw)-k_g/2\\\hline
6&4&D_4(a_1)&3&-1/2\\
7&5&D_4(a_1)&8&-3/4\\
7&5&D_4(a_1)+2A_1&8&1/4\\
7&5&2A_3+A_1&4&3/4\\
7&5&D_4(a_1)+A_1&8&-1/4
\end{array}
\]
Complex conjugation changes the displayed imbalance to its negative
and leaves the local norm class unchanged. All reflections in \(J_m\)
have norm local invariant and real entries. In rank six the sole
minimal nonreal class is represented by the particularly short word
\[
                  w=s_3s_4s_3s_2s_4s_5.
\]
Its characteristic polynomial is \((x^2+1)^2(x-1)^2\).
The rank-seven words and every corresponding moving basis and Gram
matrix are retained in the exact certificate.

The reproducible files are \texttt{check\_e6\_e7.py},
\texttt{E6\_witness\_input.json}, \texttt{E7\_witness\_input.json},
\texttt{E6\_local\_invariants.json}, and
\texttt{E7\_local\_invariants.json}. The checker shares only the
elementary rational Gaussian arithmetic helper with the G31 certificate;
it imports no numerical invariant or conclusion from that calculation.

\subsection*{D-E67-02. Every dominant sixfold source is split}

First suppose the base has genus zero. A sixfold natural multiplicity
space has total moving dimension
\[
                              D=2m+6.
\]
Period rank nine requires at least twelve branch points. Every inertia
is nonidentity, so the sum of the excesses \(k_g-1\) is
\(D-r\le2m-6\).

A generating tuple for \(J_m\) contains a nonreal entry. The
homomorphism recording whether its matrix is real or purely imaginary
has values in \(\mathbb Z/2\), and the product-one relation makes
the number of nonreal entries even. Thus it contains at least two.
For \(m=6\), each contributes excess at least three; for \(m=7\),
each contributes excess at least four. In either case these two entries
already exhaust the entire permitted excess. It follows that
\[
 r=12,\quad\text{ten entries are reflections, and the other two have}
 \quad k_g=4\ (m=6),\quad k_g=5\ (m=7).
\]
The table gives local classes three and three for the latter two entries
in rank six. Their product is nine, a Gaussian norm. In rank seven each
local class is itself a Gaussian norm. The ten reflections contribute
only norms. D-DET-01, with the phase normalization explained above,
therefore proves split signed determinant in both cases. This is a
necessary-dimension obstruction combined with an exact determinant
calculation, not a claim that these branch data have maximal period rank.

For a base of genus \(\gamma\), the same dimension equation is
\(D=2m+6-2m\gamma\). If \(\gamma=1\), then \(D=6\), so there
are at most six branch points and at most six moduli parameters.
If \(\gamma\ge2\), \(D<0\) for \(m=6,7\), which is impossible.
Thus higher-genus bases cannot restore dominance.

The conclusion also applies to a finite group whose irreducible
reflection normal-closure component is of type \(E_6\) or \(E_7\)
over \(K\). Indeed its stabilizer acts by a finite unitary normalizer
of that Weyl group. The rational-intertwiner argument of D-FIELD-02
writes such a normalizer as a scalar times a rational similarity.
All roots have the same length, so after rational rescaling the similarity
is orthogonal and permutes the roots. An automorphism of the root system
is a Weyl element followed by a diagram automorphism: move its chamber
back by the simply transitive Weyl action and inspect the simple roots.
For \(E_7\) the diagram has no automorphism, and for \(E_6\) its
involution is induced by \(-w_0\), where \(w_0\) is the longest
Weyl element. Hence the normalizer lies in scalar multiples of the
Weyl group. A scalar in \(K\) giving a finite group is in \(\mu_4\).
If the stabilizer character is nonrational it contains a nonreal such
element and, since it contains the Weyl group, it contains \(iI\);
its image is exactly \(J_m\).

Induction from the block stabilizer identifies the multiplicity
variation and actual polarization on the intermediate quotient curve,
as proved in D-FIELD-03. Applying the preceding obstruction there shows
that no such induced finite-group source dominates a nonsplit Weil-sixfold
component. This result concerns \(E_6,E_7\); it makes no assertion
that they exhaust the Gaussian reflection-normal-closure possibilities.
